\documentclass[UTF-8, a4paper,landscape]{ctexart}
\usepackage[margin=10mm]{geometry}
\usepackage{comment}

% ---------------------------------------------------------------------------
% Engine: XeLaTeX (UTF-8 native, no inputenc / CJKutf8 / CJK* environment).
% Fonts are configured in one place below; change only these lines to reskin.
% Fandol ships with TeX Live, so this compiles on any box without installing
% system fonts. Alternatives are listed in the comments.
% ---------------------------------------------------------------------------
\usepackage{fontspec}
\usepackage{xeCJK}

% Latin faces (fontspec defaults to Latin Modern; stated explicitly on purpose).
\setmainfont{Latin Modern Roman}
\setsansfont{Latin Modern Sans}
\setmonofont{Latin Modern Mono}

% CJK faces. Fandol has no true italic, so Kai stands in for it.
%   Linux alternative : \setCJKmainfont{Noto Serif CJK SC}[BoldFont=Noto Sans CJK SC Bold]
%   Windows alternative: \setCJKmainfont{SimSun}[BoldFont=SimHei, ItalicFont=KaiTi]

\begin{comment}
\setCJKmainfont{FandolSong-Regular.otf}[
  BoldFont   = FandolHei-Regular.otf,
  ItalicFont = FandolKai-Regular.otf,
]
\end{comment}
\setCJKmainfont{LXGWWenKai-Regular}
% \setCJKmainfont{HYSunShangXiang-F}
% \setCJKmainfont{FZJLJW--GB1-0}

\setCJKsansfont{FandolHei-Regular.otf}
\setCJKmonofont{FandolSong-Regular.otf}

% Route the arrow block to the CJK font, which covers U+2190..U+21FF.
\xeCJKDeclareCharClass{CJK}{"2190 -> "21FF}

\usepackage{tikz}
\usetikzlibrary{arrows.meta,positioning,calc,fit,shapes.geometric,backgrounds}
\usepackage{microtype}
\pagestyle{empty}
\setlength{\parindent}{0pt}

% ---------------------------------------------------------------------------
% Horizontal budget (A4 landscape, 10mm margins => 277mm of text width):
%   node width = 36mm text width + 2 x 2mm inner xsep = 40mm
%   row width  = 6 x 40mm + 5 x 7mm gap             = 275mm
% All nodes are placed with the "positioning" library, i.e. relative to their
% predecessor, so the gaps are controlled by a single "node distance" value.
% ---------------------------------------------------------------------------
\tikzset{
  flowstep/.style={
    rounded corners=3pt,
    draw=blue!55!black,
    line width=0.9pt,
    fill=blue!5,
    align=left,
    text width=36mm,
    minimum height=38mm,
    inner xsep=2mm,
    inner ysep=2.5mm,
    font=\small,
  },
  verify/.style={
    flowstep,
    draw=green!50!black,
    fill=green!6,
  },
  caution/.style={
    flowstep,
    draw=orange!65!black,
    fill=orange!8,
  },
  final/.style={
    flowstep,
    draw=violet!55!black,
    fill=violet!7,
  },
  basebox/.style={
    draw=red!50!black,
    fill=red!3,
    rounded corners=3pt,
    align=left,
    text width=268mm,
    minimum height=16mm,
    inner xsep=3mm,
    inner ysep=2.5mm,
    font=\scriptsize,
  },
  flow/.style={-{Stealth[length=3mm]}, line width=1.2pt, draw=gray!65!black},
  smallnote/.style={font=\scriptsize\itshape, text=gray!55!black, align=left},
}
\begin{document}

\begin{center}
{\LARGE\bfseries Windows GVim + coc.nvim 安装与验证流程}\par
\end{center}

\vspace{2mm}

\begin{center}
% node distance = <vertical> and <horizontal>
\begin{tikzpicture}[node distance=16mm and 7mm]

% Row 1: left -> right. Only s1 is placed freely; the rest follow relatively.
\node[flowstep] (s1) {\textbf{1  准备目录与 Git}\par
  确认 GVim 已安装；准备 \texttt{vimfiles/\allowbreak bundle}\par
  \texttt{git clone https://github.com/\allowbreak neoclide/\allowbreak coc.nvim.git coc.nvim}\par
  \texttt{git remote -v}};

\node[flowstep, right=of s1] (s2) {\textbf{2  固定 release 分支}\par
  \texttt{cd vimfiles/\allowbreak bundle/\allowbreak coc.nvim}\par
  \texttt{git checkout -b release origin/release}\par
  不要直接使用未构建的 \texttt{master}};

\node[verify, right=of s2] (s3) {\textbf{3  检查 GVim 基线}\par
  Vim \textgreater{}= 9.0.0438\par
  \texttt{:version}\par
  确认 \texttt{encoding=utf-8}};

\node[flowstep, right=of s3] (s4) {\textbf{4  安装 Node.js}\par
  官方安装包；生产环境优先 LTS\par
  当前建议：Node.js 24 LTS\par
  coc.nvim 要求 \textgreater{}= 20.19.0};

\node[verify, right=of s4] (s5) {\textbf{5  CMD 验证 Node / npm}\par
  \texttt{node --version}\par
  \texttt{npm --version}\par
  \texttt{where node}\par
  确认 PATH 生效};

\node[verify, right=of s5] (s6) {\textbf{6  启动 GVim，检查 Coc}\par
  重启 GVim 后执行\par
  \texttt{:CocInfo}\par
  不应出现 \texttt{javascript bundle not found}};

% Row 2: right -> left, hanging below the last node of row 1.
\node[flowstep, below=of s6] (s7) {\textbf{7  安装 Python LSP}\par
  \texttt{:CocInstall coc-pyright}\par
  \texttt{:CocList extensions}\par
  确认扩展处于 active 状态};

\node[verify, left=of s7] (s8) {\textbf{8  确认 Python 解释器}\par
  CMD：\texttt{python --version}\par
  CMD：\texttt{where python}\par
  工程建议使用项目虚拟环境};

\node[verify, left=of s8] (s9) {\textbf{9  实测自动补全}\par
  打开 \texttt{*.py}\par
  输入 \texttt{import os; os.}\par
  验证补全、诊断、跳转定义};

\node[final, left=of s9] (s10) {\textbf{10  配置 vimrc}\par
  设置 Tab / Enter 补全\par
  \texttt{updatetime=300}\par
  \texttt{nowritebackup / nobackup}\par
  按团队规范管理配置};

\node[caution, left=of s10] (s11) {\textbf{11  检查插件冲突}\par
  已有 deoplete 等自动补全时，避免重复接管 Tab\par
  可用 \texttt{:verbose imap <Tab>} 排查映射};

\node[final, left=of s11] (s12) {\textbf{12  团队可复制 / 固化}\par
  记录 Vim、Node、\par coc.nvim commit\par
  记录扩展版本 / Python 环境\par
  更新前保留 \texttt{vimrc} 与插件状态};

% Snake connectors
\draw[flow] (s1.east) -- (s2.west);
\draw[flow] (s2.east) -- (s3.west);
\draw[flow] (s3.east) -- (s4.west);
\draw[flow] (s4.east) -- (s5.west);
\draw[flow] (s5.east) -- (s6.west);
% s7 sits exactly below s6, so the turn is a straight vertical arrow.
\draw[flow] (s6.south) -- node[smallnote, midway, anchor=east, xshift=-3mm] {右侧掉头 ↓} (s7.north);
\draw[flow] (s7.west) -- (s8.east);
\draw[flow] (s8.west) -- (s9.east);
\draw[flow] (s9.west) -- (s10.east);
\draw[flow] (s10.west) -- (s11.east);
\draw[flow] (s11.west) -- (s12.east);

% Direction label for the second row, parked in the gap between the two rows.
\node[smallnote, above=2mm of s12] {第二排继续 ←};

% Security / reproducibility footer, aligned with the left edge of the grid.
\node[basebox, below=10mm of s12.south west, anchor=north west] (base)
  {\textbf{安全与可维护性基线：}
  通过 HTTPS 从官方 GitHub 获取 coc.nvim；普通环境使用 \texttt{release} 分支，只有源码开发才从 \texttt{master} 构建。\quad
  Node.js 优先使用受支持的 LTS；Coc 扩展使用 \texttt{:CocInstall}，避免不必要的全局 npm 安装。\quad
  团队环境建议额外记录 commit / 扩展版本，并在变更前备份 \texttt{vimrc}。};

% Source notes, anchored to the two bottom corners of the footer box.
\node[smallnote, below=5mm of base.south west, anchor=north west, text width=110mm]
  {依据：coc.nvim 官方文档 / GitHub；Pyright 官方安装说明；Node.js 官方版本表。\par
   2026-08-21：Vim \textgreater{}= 9.0.0438；Node.js \textgreater{}= 20.19.0；Node.js 24 为 LTS。};

\node[smallnote, below=5mm of base.south east, anchor=north east, text width=95mm, align=right]
  {Git remote：\texttt{https://github.com/\allowbreak neoclide/\allowbreak coc.nvim.git}\par
   Python 扩展：\texttt{coc-pyright}};

\end{tikzpicture}
\end{center}

\end{document}
